DGFiP, IRCOM des revenus 2024 (répartition) ; Insee, comptes nationaux 2025 (IR, CSG, CRDS) ; rapport d'évaluation « Financement » annexé au projet de loi d'approbation des comptes de la Sécurité sociale 2025 (CSG par assiette) ; calculs de l'auteur (synthese/modele\_purs.py, situation\_ircsg.py).\par Note : champ de la fusion, c'est-à-dire les revenus imposés au barème (salaires, pensions, revenus des indépendants, loyers, allocations chômage) ; l'IR des revenus du capital au prélèvement forfaitaire unique (environ 12,4 Md\euro{}) n'y est pas. Montants de 2024 (IRCOM, Sécurité sociale) portés en 2025 avec la croissance de l'IR des comptes nationaux (+7,9 \%, impôt encaissé en 2025), de la CSG sur les revenus d'activité (+2,3 \%) et de remplacement (+3,3 \%) et de la CRDS (+2,2 \%). Au niveau national (comptes nationaux 2025) : IR 103,6 Md\euro{}, CSG 156,6 Md\euro{}, CRDS 9,3 Md\euro{}. La CSG des pensions de 2024 (25,8 Md\euro{}) est un chiffre de la CCSS repris par le modèle, non relu à la source.
